\begin{scanpage}{0186}{168}
\[
\left.\begin{aligned}\overline\Sigma&=\overline\Sigma^0+\overline\Sigma^1,\\\Sigma&=\Sigma^0+\Sigma^1,
\end{aligned}\right\}\tag{26.16}
\]
where the functions $\Sigma^0$, $\overline\Sigma^0$ and $\Sigma^1$, $\overline\Sigma^1$ refer to the field radiated by the uniform and nonuniform part of the current, respectively. Substituting these expressions into Equations (26.04) and (26.06) and taking into account relationship (26.15), let us find the fringing field passing through the polarizer P toward the source Q. In the case of E-polarization, it equals
\[
\left.\begin{aligned}E_x=-H_\vartheta&=\frac{iaE_{0x}}4(\Sigma^1+\overline\Sigma^1)\frac{e^{ikR}}R e^{i\pi/2},\\
E_\vartheta=H_x&=\frac{iaE_{0x}}4(2\Sigma^0+\Sigma^1-\overline\Sigma^1)\frac{e^{ikR}}R,
\end{aligned}\right\}\tag{26.17}
\]
and in the case of H-polarization
\[
\left.\begin{aligned}H_x=E_\vartheta&=\frac{iaH_{0x}}4(\Sigma^1+\overline\Sigma^1)\frac{e^{ikR}}R e^{i\pi/2},\\
H_\vartheta=-E_x&=\frac{iaH_{0x}}4(2\Sigma^0+\Sigma^1-\overline\Sigma^1)\frac{e^{ikR}}R.
\end{aligned}\right\}\tag{26.18}
\]
The physical meaning of the result obtained is as follows. The field scattered by the body at the point Q is the sum of two waves polarized in mutually perpendicular directions. The reflected wave which is polarized the same as the primary radiation of the source is determined by the function $\Sigma_+=\frac12(\Sigma^1+\overline\Sigma^1)$, and is created only by the nonuniform part of the current. The reflected wave with the perpendicular polarization is described by the function $\Sigma_-=2\Sigma^0+\Sigma^1-\overline\Sigma^1$, and is the field radiated by both parts of the current. Let us note that in the general case the functions $\Sigma^1$ and $\overline\Sigma^1$ do not coincide, and therefore they are not balanced out in the expressions for $\Sigma_-$. In other words, the field radiated by the uniform part of the current in this case may not be separated from the fringing field.

Thus, the investigated method allows one to separate from the total field scattered by any metal body of finite dimensions that part of the field which is caused by a distortion of the surface (the
\end{scanpage}
